\chapter{绪论}

\section{选题背景、目的和意义}

当前制造业企业正处于由单点信息化向全流程数字化协同转型的关键阶段。订单获取、库存控制、生产执行、原材料采购以及质量管理等环节相互耦合，任何单一节点的信息延迟或处理失误，都会沿供应链链路放大并传导至交付周期、库存成本和客户满意度等核心指标。尤其是在多角色共同参与的企业业务场景下，如果订单状态、库存数量、采购进度和质检结果无法实现统一管理与实时共享，则容易形成信息孤岛，造成重复录入、人工催办、库存积压和追溯困难等问题。因此，构建一套面向企业内部多部门协同和上下游业务联动的供应链协同管理系统，已成为提升企业运营效率和管理精细化水平的重要途径\citep{christopher2016,lee2004triplea}。

从企业管理视角看，供应链管理系统已不再局限于订单记录和库存统计，而是逐步发展为集订单驱动、库存联动、生产计划、采购协同、质量追溯和消息通知于一体的综合业务平台。系统建设的重点不只是实现基础数据电子化，还在于建立跨岗位、跨流程、跨状态的协同机制，使顾客、销售管理员、仓库管理员、生产管理员、采购管理员、供应商、质检管理员和系统管理员等角色能够围绕统一数据源开展工作，并通过可控的权限边界实现职责分离与业务闭环\citep{gunasekaran2004framework,monk2012erp}。

基于上述背景，本文以企业供应链协同管理为研究对象，设计并实现了一套基于 Spring Boot 与 Vue 3 的企业供应链协同管理系统。系统采用前后端分离架构，在后端使用 Spring Boot 构建 RESTful 服务，通过 Spring Security 与 JWT 实现用户认证与授权控制，采用 MySQL 与 JPA 完成业务数据管理，并利用 WebSocket/STOMP 支持关键业务节点的实时消息推送。本文的研究意义主要体现在以下三个方面：其一，从业务层面提升订单、库存、生产、采购和质检之间的流程衔接效率；其二，从系统层面增强多角色并发访问环境下的数据一致性、可追溯性和可扩展性；其三，从工程实践层面为中小型制造企业构建轻量化、易维护的供应链协同平台提供可参考的设计思路。

\section{国内外研究现状}

\subsection{国外研究现状}

国外关于供应链管理的研究起步较早，研究内容已从单纯的物流优化扩展到面向供应链韧性、协同决策、库存配置和质量追踪的系统化管理框架。Christopher 从物流与供应链管理一体化视角强调了企业通过信息共享和流程整合提升响应能力的重要性\citep{christopher2016}。Lee 提出的“三A供应链”思想进一步指出，敏捷性、适应性与协同性是现代供应链系统建设的重要目标\citep{lee2004triplea}。在企业信息系统层面，ERP、MES 与 SCM 等系统逐渐形成互补关系，其中 ERP 更强调资源计划与财务控制，MES 更关注车间执行，而 SCM 系统则更注重跨部门和跨主体的业务协同。随着 Web 技术和服务架构的发展，面向浏览器访问、支持分布式协同的业务系统已成为主流实现方式。

在系统实现方法方面，国外研究和工业实践较多采用分层架构、无状态认证、消息驱动协同和服务接口标准化等手段提升系统的可维护性与可扩展性。针对质量追溯场景，国外企业普遍重视关键业务节点数据留痕，通过批次标识、状态记录和操作日志实现问题定位与责任界定。这些研究和实践为本文开展企业供应链协同管理系统的设计与实现提供了重要参考。

\subsection{国内研究现状}

国内在企业信息化和制造业数字化转型推动下，围绕供应链管理、库存控制、采购协同、质量追溯和制造执行的研究不断深化。现有研究普遍认为，传统企业在信息系统建设过程中存在系统割裂、角色边界不清晰、业务状态难以联动以及质量追溯链条不完整等问题，因此需要面向业务闭环重新整合订单、仓储、生产、采购和质检数据。与此同时，随着国家标准与行业规范不断完善，质量追溯体系建设受到越来越多关注，可追溯性已成为企业提升质量透明度和客户信任度的重要支撑\citep{gbt22005-2009}。

从实现路径看，国内相关系统研究更加注重工程可行性与业务适配性。一类研究侧重于 ERP、WMS、MES 等现有系统的功能扩展与对接，另一类研究则采用 Spring Boot、Vue 等主流 Web 技术重构业务流程，以实现更灵活的前后端分离与角色化访问控制。总体而言，现有研究为供应链协同系统建设提供了较好基础，但针对多角色权限隔离、订单驱动下的仓储与生产联动、库存预警触发采购协同、质检留痕与质量追溯、实时消息同步等方面，仍需结合具体业务流程开展更细致的工程设计与实现分析。

\section{研究内容与论文结构安排}

\subsection{研究内容}

围绕企业供应链协同场景，本文的主要研究内容包括以下几个方面。
\begin{enumerate}
  \item 分析系统所面向的业务环境与用户角色，明确顾客、销售管理员、仓库管理员、生产管理员、采购管理员、供应商、质检管理员和系统管理员在业务协同中的职责边界与功能需求。
  \item 设计系统总体架构、前后端分层方式、权限控制机制以及数据库结构，构建面向订单驱动的业务闭环，实现订单、库存、生产、采购、质检与追溯数据的统一管理。
  \item 实现用户认证与账号管理、顾客下单与订单管理、销售审核、仓库管理与库存预警、生产管理、采购与供应商管理、质检与质量追溯、实时消息通知以及周计划生成等核心模块。
  \item 重点研究基于角色的访问控制、订单至质检协同流程、采购供应协同流程、库存预警与周计划联动机制、实时消息推送以及质量追溯的工程实现方法，并通过系统测试验证其可用性与合理性。
\end{enumerate}

\subsection{论文结构安排}

本文共分为七章，各章内容安排如下。
\begin{enumerate}
  \item 第1章为绪论，主要阐述课题背景、研究目的与意义，总结国内外相关研究现状，并说明本文的研究内容与整体结构安排。
  \item 第2章为理论基础与相关技术，介绍 B/S 架构、前后端分离模式、Spring Boot、Spring Security、JWT、Vue 3、Vite、Axios、WebSocket/STOMP、MySQL 与 JPA 等关键技术，并分析其在本系统中的适用性。
  \item 第3章为系统分析与设计，从需求分析、总体架构、权限控制、业务流程以及数据库设计等方面完成系统的整体设计。
  \item 第4章为系统详细设计与实现，围绕各业务模块的实现过程，分析前端页面、后端接口、业务服务与数据实体之间的协同关系。
  \item 第5章为核心功能实现，重点讨论基于角色的访问控制、业务闭环协同、库存预警与周计划生成、实时消息同步以及质量追溯等关键实现机制。
  \item 第6章为系统测试，介绍测试环境、单元测试与功能测试设计，并对测试结果进行分析。
  \item 第7章为总结与展望，对全文工作进行总结，分析系统存在的不足，并提出后续优化方向。
\end{enumerate}